\chapter*{ประมวลรายวิชา (Course Syllabus)}
\addcontentsline{toc}{chapter}{ประมวลรายวิชา (Course Syllabus)}

\section*{ข้อมูลทั่วไปของรายวิชา}
\begin{tabbing}
\hspace{4.0cm}\=\kill
\textbf{รหัสและชื่อรายวิชา} \> AGR2105 จุลชีววิทยาสำหรับการเกษตร (Microbiology for Agriculture)\\
\textbf{จำนวนหน่วยกิต} \> 3(2-2-5) (บรรยาย 2 คาบ ปฏิบัติการ 2 คาบ ศึกษาค้นคว้าด้วยตนเอง 5 ชั่วโมง)\\
\textbf{หลักสูตรและประเภท} \> วิทยาศาสตรบัณฑิต สาขาวิชาเกษตรศาสตร์ (หมวดวิชาเฉพาะด้าน)\\
\textbf{อาจารย์ผู้รับผิดชอบ} \> ผู้ช่วยศาสตราจารย์ ดร.จิรภัทร จันทมาลี\\
\textbf{สถาบันการศึกษา} \> คณะเทคโนโลยีการเกษตร มหาวิทยาลัยราชภัฏรำไพพรรณี
\end{tabbing}

\section*{คำอธิบายรายวิชา (Course Description)}
บทบาทและความสำคัญของจุลินทรีย์ในระบบนิเวศการเกษตร อนุกรมวิธานและความหลากหลายของจุลินทรีย์ในดิน วัฏจักรชีวธรณีเคมีของธาตุอาหารพืช ปฏิสัมพันธ์ระหว่างจุลินทรีย์กับพืชบริเวณไรโซสเฟียร์ แบคทีเรียส่งเสริมการเจริญเติบโตของพืช (PGPR) เชื้อราไมคอร์ไรซา การตรึงไนโตรเจนทางชีวภาพ การควบคุมศัตรูพืชโดยชีววิธี จุลินทรีย์ก่อโรคพืชและการจัดการ การหมักของเสียทางการเกษตร การผลิตปุ๋ยชีวภาพและก๊าซชีวภาพ เทคโนโลยีชีวภาพทางจุลชีววิทยาการเกษตร บทบาทของจุลินทรีย์ต่อการเปลี่ยนแปลงสภาพภูมิอากาศ และการประยุกต์ใช้เทคโนโลยีโอมิกส์สมัยใหม่ร่วมกับปัญญาประดิษฐ์ในการเกษตรแม่นยำ

\section*{ผลลัพธ์การเรียนรู้ระดับรายวิชา (Course Learning Outcomes: CLOs)}
เมื่อสิ้นสุดการเรียนการสอนรายวิชานี้ นักศึกษาสามารถ:
\begin{enumerate}[label=\textbf{CLO\arabic*:}]
    \item \textbf{อธิบายและจำแนก} ชนิด โครงสร้าง หน้าที่ และความหลากหลายของจุลินทรีย์ในดินและระบบนิเวศการเกษตรได้อย่างถูกต้องตามหลักอนุกรมวิธาน (\textit{Cognitive Level: Remember \& Understand})
    \item \textbf{วิเคราะห์} กลไกทางชีวเคมีของจุลินทรีย์ในวัฏจักรการหมุนเวียนธาตุอาหาร และปฏิสัมพันธ์ระหว่างจุลินทรีย์กับรากพืชได้อย่างเป็นระบบ (\textit{Cognitive Level: Analyze})
    \item \textbf{ประยุกต์ใช้} จุลินทรีย์ปฏิปักษ์ในการควบคุมศัตรูพืชโดยชีววิธี และออกแบบแผนการจัดการโรคพืชอย่างบูรณาการเพื่อลดการใช้สารเคมี (\textit{Cognitive Level: Apply})
    \item \textbf{ปฏิบัติการทดลอง} การเพาะเลี้ยง การคัดแยก การตรวจสอบคุณสมบัติ และการเตรียมหัวเชื้อจุลินทรีย์ทางการเกษตรได้อย่างถูกต้องและปลอดภัยตามมาตรฐานสากล (\textit{Psychomotor Domain})
    \item \textbf{บูรณาการองค์ความรู้} เทคโนโลยีดีเอ็นเอรหัสแท่ง เมตาเจโนมิกส์ และการวิเคราะห์ข้อมูลด้วยระบบอัจฉริยะ เพื่อแก้ปัญหาการจัดการดินและยกระดับผลผลิตพืชอย่างยั่งยืน (\textit{Cognitive Level: Create})
\end{enumerate}

\section*{แผนที่แสดงการกระจายความรับผิดชอบมาตรฐานผลการเรียนรู้ (Curriculum Mapping)}
\begin{table}[H]
\centering
\small
\begin{tabularx}{\textwidth}{|X|c|c|c|c|c|}
\hline
\rowcolor{agriMint}
\textbf{หัวข้อเนื้อหาประจำสัปดาห์} & \textbf{CLO1} & \textbf{CLO2} & \textbf{CLO3} & \textbf{CLO4} & \textbf{CLO5} \\ \hline
1. บทนำและประวัติศาสตร์จุลชีววิทยาการเกษตร & $\bullet$ & $\circ$ & & & \\ \hline
2. ความหลากหลายและโครงสร้างของจุลินทรีย์ในดิน & $\bullet$ & $\bullet$ & & $\bullet$ & \\ \hline
3. วัฏจักรชีวธรณีเคมีและการหมุนเวียนธาตุอาหาร & & $\bullet$ & & $\bullet$ & \\ \hline
4. จุลินทรีย์บริเวณรากพืชและการส่งเสริมการเจริญเติบโต & $\bullet$ & $\bullet$ & & $\bullet$ & $\circ$ \\ \hline
5. การตรึงไนโตรเจนทางชีวภาพและปุ๋ยชีวภาพ & & $\bullet$ & & $\bullet$ & $\bullet$ \\ \hline
6. จุลินทรีย์กับการควบคุมศัตรูพืชโดยชีววิธี & & & $\bullet$ & $\bullet$ & $\bullet$ \\ \hline
7. โรคพืชที่เกิดจากจุลินทรีย์และการจัดการแบบบูรณาการ & $\bullet$ & & $\bullet$ & $\bullet$ & \\ \hline
8. จุลินทรีย์ในการหมักและการจัดการของเสียเกษตร & & $\bullet$ & $\bullet$ & $\bullet$ & \\ \hline
9. อุตสาหกรรมเกษตรแปรรูปและเทคโนโลยีชีวภาพ & $\bullet$ & $\bullet$ & & $\bullet$ & $\circ$ \\ \hline
10. จุลินทรีย์กับสภาพภูมิอากาศและการกักเก็บคาร์บอน & & $\bullet$ & & & $\bullet$ \\ \hline
11. การประยุกต์ใช้ AI และเมตาเจโนมิกส์ในดินเกษตร & & $\circ$ & & $\bullet$ & $\bullet$ \\ \hline
\end{tabularx}
\caption*{สัญลักษณ์: $\bullet$ ความรับผิดชอบหลัก \quad $\circ$ ความรับผิดชอบรอง}
\end{table}

\clearpage
